\section{The Persistence Enclosure Theorem}
\label{sec:theorem}

This section proves Theorem~\ref{thm:persistence-enclosure} by exhaustive
case analysis. The structure mirrors the Constitutional Enclosure proof
of BTV~\cite{btv2026} (Section~4.4), extending it with a fourth case
for the cryptographic boundary introduced by the network-adjacent adversary.

\subsection{The Enclosure Perimeter}

The type law (\ref{eq:delivery-law}) defines a \emph{perimeter}:
the set of code paths that could, in principle, allow a
\texttt{DeliveryToken} to be constructed without a valid
\texttt{InclusionReceipt}. We enumerate every such path and show
each is closed by a distinct, compiler-enforced barrier.

\begin{theorem}[Persistence Enclosure]
\label{thm:pe-full}
In any Safe Rust program $P$ whose crate graph includes
\texttt{btv-transparency}, the set of Ephemeral Verdicts
$\mathcal{E}(P) = \emptyset$.
Equivalently, if $P$ type-checks under Safe Rust, then every
delivered verdict is backed by a valid \texttt{InclusionReceipt}.
\end{theorem}

\begin{proof}
We identify four exhaustive attack cases and show each is
ill-typed or cryptographically rejected before the receipt enters
the linear pipeline.

\paragraph{Case A: Struct-Literal Forgery (E0451).}
An adversary constructs \texttt{DeliveryToken~\{~verdict\_bytes:~\ldots,
log\_index:~\ldots,~signature\_bytes:~\ldots~\}} directly.
\begin{itemize}
  \item \emph{Defence.} All fields of \texttt{DeliveryToken} are declared
  without \texttt{pub} visibility; they are private to the
  \texttt{btv-transparency} crate.
  \item \emph{Compiler error.} \texttt{error[E0451]: field
  `verdict\_bytes` of struct `DeliveryToken` is private}.
  \item \emph{Verified by.} Compile-fail test
  \texttt{tests/ui/delivery\_struct\_literal.rs}.
\end{itemize}

\paragraph{Case B: Receipt Forgery via External Constructor (E0603).}
An adversary calls \texttt{InclusionReceipt::new(index,~sig,~key)}
directly to manufacture a receipt without contacting the Log.
\begin{itemize}
  \item \emph{Defence.} \texttt{InclusionReceipt::new} is declared
  \texttt{pub(crate)}; it is not part of the public API surface.
  The sole public path to an \texttt{InclusionReceipt} is
  \texttt{LogClient::submit\_and\_await}, which contacts the Log.
  \item \emph{Compiler error.} \texttt{error[E0603]: function
  `new` is private}.
  \item \emph{Verified by.} Compile-fail test
  \texttt{tests/ui/delivery\_without\_receipt.rs}.
\end{itemize}

\paragraph{Case C: Silent Drop of Receipt (\texttt{unused\_must\_use}).}
A developer calls \texttt{LogClient::submit\_and\_await} but discards
the result (\texttt{let~\_~=~client.submit\_and\_await(\&hash);})
and proceeds to call \texttt{DeliveryToken::seal} via an alternative
path.
\begin{itemize}
  \item \emph{Defence.} \texttt{InclusionReceipt} carries
  \texttt{\#[must\_use = "\ldots"]}. The crate-level attribute
  \texttt{\#![deny(unused\_must\_use)]} promotes the warning to a
  hard compilation error. Furthermore, no alternative path to
  \texttt{DeliveryToken::seal} exists without a value of type
  \texttt{InclusionReceipt} (closed by Case~B).
  \item \emph{Compiler error.} \texttt{error: unused `InclusionReceipt`
  which must be used}.
  \item \emph{Verified by.} Compile-fail test
  \texttt{tests/ui/receipt\_drop\_silent.rs}.
\end{itemize}

\paragraph{Case D: Network MITM / Forged Signature.}
A System Operator intercepts the HTTP response from the Log and
substitutes a receipt signed with a key under the operator's control,
bypassing the Log authority assumption at the network layer.
\begin{itemize}
  \item \emph{Defence.} \texttt{LogClient} is constructed with a
  \texttt{VerifyingKey} obtained out-of-band (e.g., pinned at build
  time or from a trusted registry). Before constructing the
  \texttt{InclusionReceipt}, \texttt{submit\_and\_await} verifies the
  Ed25519 signature against this pinned key. A forged or substituted
  response fails verification and returns
  \texttt{Err(LogClientError::InvalidSignature)}; no receipt is
  produced.
  \item \emph{Consequence.} By Case~B, the caller has no alternative
  path to an \texttt{InclusionReceipt}. The pipeline fails securely
  (Corollary~\ref{cor:fail-secure}).
  \item \emph{Verified by.} Unit test
  \texttt{invalid\_signature\_rejected} in \texttt{src/lib.rs}.
\end{itemize}

\medskip
Cases A--D are exhaustive: they cover every path by which an
\texttt{InclusionReceipt} could enter the pipeline without a genuine
Log submission. Since \texttt{DeliveryToken::seal} requires a value
of type \texttt{InclusionReceipt} and all forgery paths are closed,
no \texttt{DeliveryToken} can be constructed without a Log-issued
receipt. Since \texttt{DeliveryToken::deliver} is the sole function
that crosses the API boundary (its signature consumes \texttt{self}
by value), no delivery can occur without a valid receipt. Therefore
$\mathcal{E}(P) = \emptyset$. $\square$
\end{proof}

\subsection{Corollary: Fail-Secure Persistence}

\begin{corollary}[Fail-Secure Persistence]
\label{cor:fail-secure}
If the Transparency Log is unavailable or returns an error,
\texttt{LogClient::submit\_and\_await} returns \texttt{Err}.
Because \texttt{DeliveryToken::seal} requires an
\texttt{InclusionReceipt} by value (not \texttt{Option<InclusionReceipt>}
or \texttt{Result<\ldots>}), the error-handling code has no
well-typed path to a \texttt{DeliveryToken}. The system therefore
cannot deliver a decision when the log is unreachable.
This extends Corollary~4.7 of BTV~\cite{btv2026} (fail-secure
evidence) to the persistence layer.
\end{corollary}

\begin{proof}
By the type of \texttt{seal}: \texttt{fn~seal(verdict:~\&impl~Serialize,
receipt:~InclusionReceipt)~->~Result<DeliveryToken,~SealError>}.
The second argument is \texttt{InclusionReceipt}, not a nullable or
fallible wrapper. A caller that handles the \texttt{Err} branch of
\texttt{submit\_and\_await} does not possess an
\texttt{InclusionReceipt} value in scope. Cases~A and~B close all
forgery paths. Therefore no well-typed term of type
\texttt{DeliveryToken} can be constructed in the error branch.
$\square$
\end{proof}

\subsection{Relationship to the Constitutional Enclosure Theorem}

Table~\ref{tab:comparison} compares the two enclosure theorems
structurally. Both proofs share the same three-case compile-time
perimeter (Cases~A--C). BTV-Transparency adds Case~D, which bridges
the type-system boundary to the cryptographic infrastructure layer.
The two theorems compose: a program satisfying both invariants
cannot produce a decision that (a) lacks an evidence chain, or
(b) was delivered without a verifiable persistence record.

\begin{table}[h]
\caption{Structural comparison of enclosure theorems}
\label{tab:comparison}
\centering
\begin{tabular}{llll}
\hline
\textbf{Case} & \textbf{Attack} & \textbf{BTV [1]} & \textbf{BTV-T (this paper)} \\
\hline
A & Struct-literal forge & E0451 & E0451 \\
B & Private constructor & E0603/E0624 & E0603 \\
C & Silent drop & must\_use & must\_use \\
D & Crypto MITM & --- & Ed25519 verify \\
\hline
\end{tabular}
\end{table}
